\section{Introducción}

La dispersión de la luz por partículas es un problema clásico de la
electrodinámica que sigue siendo relevante en muchos contextos actuales. Cuando
una onda electromagnética incide sobre una partícula, el campo induce en ella
cargas y corrientes de polarización que reemiten parte de la energía incidente
en distintas direcciones. Cómo se reparte esa energía depende del tamaño, la
geometría y la composición de la partícula, y de la longitud de onda de la luz
incidente. El primer tratamiento completo de este problema para una esfera lo
dieron Ludvig Lorenz y, de forma independiente y más conocida, el físico alemán
Gustav Mie en 1908. La teoría de Mie proporciona la solución analítica exacta
para una esfera homogénea, a partir de la cual se calculan las
secciones eficaces de dispersión, absorción y extinción, así como la
distribución angular de la radiación dispersada. Más de un siglo después, sigue
siendo una herramienta de referencia para entender la interacción entre la luz y
la materia.

La esfera aislada es solo el punto de partida. Muchos sistemas de interés no
están formados por una única partícula, sino por conjuntos de esferas próximas
entre sí, donde el campo que excita a cada una no es únicamente el incidente
externo, sino también el dispersado por las demás. Aparece entonces un problema
de dispersión múltiple, en el que las respuestas individuales quedan acopladas.
Este caso se aborda extendiendo la propia teoría de Mie: cada esfera se sigue
describiendo mediante sus coeficientes de Mie, pero acoplados entre sí a través
del formalismo de la matriz de transición, o matriz $T$, que relaciona los
coeficientes del campo excitante con los del campo dispersado en una base de
armónicos esféricos vectoriales. Este enfoque está implementado en el código
MSTM (\textit{Multiple Sphere T-Matrix}), escrito en Fortran por Mackowski, que
constituye la herramienta de cálculo sobre la que se apoya este trabajo y que
permite estudiar tanto propiedades de campo lejano como distribuciones de campo
cercano.

La dispersión por partículas esféricas aparece en disciplinas muy distintas. En
óptica atmosférica permite comprender fenómenos asociados a aerosoles, gotas y
partículas en suspensión; en física biomédica interviene en la interacción de la
luz con tejidos y partículas biológicas; y en nanofotónica resulta esencial para
diseñar estructuras capaces de controlar la luz por debajo de la escala de la
longitud de onda. Dentro de este último campo, las nanopartículas metálicas
---y en particular las de oro--- son capaces de concentrar la luz incidente en
regiones mucho más pequeñas que su longitud de onda, intensificando
notablemente el campo electromagnético en su entorno inmediato. Estas regiones
de campo intenso, o \emph{hot spots}, son la base de las nanoantenas ópticas y
sus aplicaciones. Aunque este trabajo no desarrolla una aplicación experimental concreta, sí
estudia los mecanismos físicos que las sustentan.

El objetivo de este trabajo es estudiar la dispersión electromagnética por
nanopartículas esféricas de oro y por agregados simples de esferas mediante
simulaciones numéricas basadas en MSTM. Para ello se toma primero la esfera
individual como referencia y después se analizan configuraciones
acopladas, concretamente un dímero y un tetrámero. 
El estudio considera tanto la respuesta en campo lejano como la intensificación
del campo cercano alrededor de las partículas. Finalmente, los resultados
obtenidos se reinterpretan desde el punto de vista de las nanoantenas ópticas,
analizando el papel del tamaño, la separación entre partículas y la geometría
del agregado.

La utilización de MSTM implica preparar los ficheros de entrada en un formato
concreto, trabajar en unidades adimensionales y procesar los ficheros de salida
que genera el programa. Para agilizar este proceso se ha desarrollado un entorno
gráfico propio en Python, que automatiza la generación de los ficheros de entrada, la
ejecución de MSTM y el postproceso de los resultados.

A partir de este planteamiento, los objetivos específicos del trabajo son:

\begin{enumerate}
\item Comprender el formalismo de Mie y de la matriz $T$, así como su aplicación
        al problema de la dispersión electromagnética por una o varias esferas.

\item Analizar la implementación de estos métodos en el código MSTM de
        Mackowski, identificando las magnitudes físicas que calcula y la forma
        en que deben interpretarse sus resultados.

\item Desarrollar un entorno de trabajo en Python que automatice la
        configuración, ejecución y postproceso de simulaciones con MSTM.

\item Estudiar la dispersión de la luz por nanopartículas esféricas de oro,
        tomando la esfera individual como referencia y analizando el efecto del
        acoplamiento electromagnético en un dímero y en un tetrámero.

\item Reinterpretar los resultados obtenidos en el marco de las nanoantenas
        ópticas y sus aplicaciones.
\end{enumerate}